\abstracteng{
Understanding the neural mechanisms underlying decision making is a central goal of neuroscience. In multi‐alternative tasks, animals often display idiosyncratic choice biases—consistent preferences not explained by external cues—which may arise from subtle asymmetries in their underlying circuits. Here, we develop a mechanistic rate‐based network model to explain how such biases emerge and how temporal features of the task influence performance in a three‐choice delayed‐response paradigm (3CDR). The model comprises three excitatory populations (one per choice) coupled via a global inhibitory population and driven by both a non‐selective “urgency” ramp and transient, location‐specific inputs. By performing a normal‐form reduction, we map its dynamics onto a two‐dimensional drift‐diffusion process that is equivalent to the motion of a particle on an evolving energy landscape. Stochastic simulations quantitatively reproduce mice choice accuracy across variations in stimulus and delay durations. Our framework thus provides a quantitative tool for interpreting neural recordings and for designing targeted perturbations to probe how circuit asymmetries and timing signals shape decision outcomes.}